
定义
\begin{equation}
\begin{aligned}
\widetilde{\mathfrak C}_X
={}&C_1([0,1]^2\times L'_1,
(\lambda'_0,\lambda'_1))\vee_r{}\\
&C_2([0,1]^2\times L'_1,
(\lambda'_0,\lambda'_1)).
\end{aligned}
\tag{7.2.353}\label{eq:7-2-353}
\end{equation}
取 $[0,1]^2\times C(L'_1,\lambda'_1)$
的 $A_{n',K'}$ 类模型
$\widetilde{\mathfrak C}_4$，满足
\[
\begin{aligned}
\widetilde{\mathfrak C}_{4,s=s_0}
&=[0,1]_r\times C(L'_1,\lambda'_1),\\
\widetilde{\mathfrak C}_{4,r=r_0}
&=\mathfrak C'_1
=C([0,1]_s\times L'_1,\lambda'_1).
\end{aligned}
\]
后一端面条件即 (7.2.278.1)。
这里略微滥用记号：
$[0,1]_r\times C(L'_1,\lambda'_1)$
表示命题~\ref{prop:7-2-190}
所构造的任意一个此类模型。
由定理~\ref{thm:4-2-34}，存在
\[
\operatorname{id}\times j\colon
(\widetilde{\mathfrak C}_X)_{s=s_0}
\longrightarrow
[0,1]_r\times C(L'_1,\lambda'_1),
\]
对构成胶合的 $i=1,2$ 两块均如此
\translatornote{源书在此写“for $i=1,2$”，
但显示的 $\operatorname{id}\times j$ 及下式均无 $i$ 下标；
按 (7.2.353) 将其理解为两块上的相容构造，
仍沿用源书的无下标记号。}；
它与 $\operatorname{Incl}$ 交换，且
\begin{equation}
\operatorname{Eval}_{r=r_0}
\circ(\operatorname{id}\times j)
=j\colon C(L'_1,(\lambda'_0,\lambda'_1))
\longrightarrow C(L'_1,\lambda'_1).
\tag{7.2.354}\label{eq:7-2-354}
\end{equation}

\begin{lemma}{7.2.355}\label{lem:7-2-355}
存在滤过 $A_{n',K'}$ 同态
\[
\widetilde{\mathfrak I}\colon
\widetilde{\mathfrak C}_X
\longrightarrow\widetilde{\mathfrak C}_4,
\]
满足下列性质：
\begin{equation}
\widetilde{\mathfrak I}\circ\operatorname{Incl}
=\operatorname{Incl},
\tag{7.2.356.1}\label{eq:7-2-356-1}
\end{equation}
\begin{equation}
\operatorname{Eval}_{r=r_0}\circ
\widetilde{\mathfrak I}
=j_{[0,1]_s}\circ
\operatorname{Eval}_{r=r_0},
\tag{7.2.356.2}\label{eq:7-2-356-2}
\end{equation}
\begin{equation}
\operatorname{Eval}_{s=s_0}\circ
\widetilde{\mathfrak I}
=(\operatorname{id}\times j)
\circ\operatorname{Eval}_{s=s_0}.
\tag{7.2.356.3}\label{eq:7-2-356-3}
\end{equation}
\end{lemma}

\noindent\textbf{证明。}
利用命题~\ref{prop:7-2-312}，
可由纯代数方法直接证明。
定义
\begin{equation}
\widetilde{\mathfrak C}
=\widetilde{\mathfrak C}_4
\vee_r\widetilde{\mathfrak C}_3,
\tag{7.2.357}\label{eq:7-2-357}
\end{equation}
\begin{equation}
\begin{aligned}
\widetilde{\mathfrak H}
={}&\bigl(\widetilde{\mathfrak I}\circ
(\widetilde{\mathfrak H}_1
\vee_r\widetilde{\mathfrak H}_2)\bigr)\\
&\vee_r\bigl(\widetilde{\mathfrak H}_3
\circ\mathfrak f^{\phi'}\bigr).
\end{aligned}
\tag{7.2.358}\label{eq:7-2-358}
\end{equation}
容易验证，它们满足所需性质。\qed

我们指出，本节已详细讨论滤过 $A_\infty$ 代数、
滤过 $A_\infty$ 同态以及它们之间的
滤过 $A_\infty$ 同伦的构造。
这些证明的写法，使它们能够直接推广到
滤过 $A_\infty$ 双模、滤过 $A_\infty$ 双模同态
及其之间的滤过 $A_\infty$ 同伦。
同伦含幺版本的推广将在下一节讨论。
第 7.4 节还会进一步推广，以处理算子
$\mathfrak p$、$\mathfrak q$、$\mathfrak r$。
